\documentclass[11pt, a4paper]{report}

% ========== Common Packages ==========
\usepackage{fontspec}
\usepackage{kotex}
\usepackage{amsmath, amssymb, amsthm}
\usepackage{mathtools}
\usepackage{bm}
\usepackage{graphicx}
\usepackage{booktabs}
\usepackage{array}
\usepackage{tabularx}
\usepackage{longtable}
\usepackage{hyperref}
\usepackage{cleveref}
\usepackage{geometry}
\usepackage{fancyhdr}
\usepackage{titlesec}
\usepackage{enumitem}
\usepackage{xcolor}
\usepackage{tcolorbox}
\usepackage{algorithm}
\usepackage{algorithmic}
\usepackage{tikz}
\usetikzlibrary{arrows.meta, positioning, calc, decorations.pathreplacing, shapes.geometric, fit, patterns, angles, quotes, trees}
\usepackage{pgfplots}
\pgfplotsset{compat=1.18}
\usepackage{caption}
\usepackage{subcaption}
\usepackage{float}

\geometry{margin=2.5cm, top=3cm, bottom=3cm}

\usepackage{multirow}
% ========== Colors ==========
% Common
\definecolor{defblue}{RGB}{30, 80, 150}
\definecolor{thmgreen}{RGB}{20, 110, 60}
\definecolor{noteorange}{RGB}{200, 100, 0}
\definecolor{lightblue}{RGB}{230, 240, 255}
\definecolor{lightgreen}{RGB}{230, 255, 235}
\definecolor{lightorange}{RGB}{255, 245, 230}
\definecolor{lightgray}{RGB}{245, 245, 245}
% Extended
\definecolor{lightred}{RGB}{255, 235, 235}
\definecolor{darkred}{RGB}{180, 30, 30}
\definecolor{biogreen}{RGB}{10, 130, 80}
\definecolor{cvpurple}{RGB}{100, 40, 150}
\definecolor{injuryred}{RGB}{200, 40, 40}
\definecolor{codegray}{RGB}{240, 240, 245}
\definecolor{pipeblue}{RGB}{25, 100, 180}
\definecolor{constraintred}{RGB}{200, 50, 50}
\definecolor{termblack}{RGB}{30, 30, 30}
\definecolor{goldcolor}{RGB}{180, 140, 20}


% ========== Korean theorem names (이 파일 전용) ==========
\theoremstyle{definition}
\newtheorem{definition}{정의}[chapter]
\newtheorem{theorem}{정리}[chapter]
\newtheorem{proposition}{명제}[chapter]
\newtheorem{corollary}{따름정리}[chapter]
\newtheorem{example}{예제}[chapter]
\theoremstyle{remark}
\newtheorem{remark}{참고}[chapter]

% ========== Korean tcolorbox (이 파일 전용) ==========
\newtcolorbox{keyidea}[1][]{colback=lightblue, colframe=defblue, fonttitle=\bfseries, title={핵심}, #1}
\newtcolorbox{importantnote}[1][]{colback=lightorange, colframe=noteorange, fonttitle=\bfseries, title={중요}, #1}
\newtcolorbox{warningbox}[1][]{colback=lightred, colframe=darkred, fonttitle=\bfseries, title={주의}, #1}
\newtcolorbox{proofbox}[1][]{colback=lightgreen, colframe=thmgreen, fonttitle=\bfseries, title={증명/유도}, #1}
\newtcolorbox{practicebox}[1][]{colback=lightgray, colframe=gray!60, fonttitle=\bfseries, title={실전 적용}, #1}

% ========== Custom math commands ==========
% Vectors (lowercase)
\newcommand{\vx}{\mathbf{x}}
\newcommand{\vd}{\mathbf{d}}
\newcommand{\vo}{\mathbf{o}}
\newcommand{\vc}{\mathbf{c}}
\newcommand{\vr}{\mathbf{r}}
\newcommand{\vp}{\mathbf{p}}
\newcommand{\vv}{\mathbf{v}}
\newcommand{\vn}{\mathbf{n}}
\newcommand{\vu}{\mathbf{u}}
\newcommand{\vt}{\mathbf{t}}
\newcommand{\ve}{\mathbf{e}}
\newcommand{\vq}{\mathbf{q}}
% Vectors (uppercase)
\newcommand{\vF}{\mathbf{F}}
\newcommand{\vM}{\mathbf{M}}
\newcommand{\va}{\mathbf{a}}
\newcommand{\vg}{\mathbf{g}}
\newcommand{\vX}{\mathbf{X}}
% Bold Greek
\newcommand{\vmu}{\bm{\mu}}
\newcommand{\vbeta}{\bm{\beta}}
\newcommand{\vtheta}{\bm{\theta}}
\newcommand{\vpsi}{\bm{\psi}}
% Matrices
\newcommand{\mSigma}{\bm{\Sigma}}
\newcommand{\mR}{\mathbf{R}}
\newcommand{\mS}{\mathbf{S}}
\newcommand{\mW}{\mathbf{W}}
\newcommand{\mJ}{\mathbf{J}}
\newcommand{\mG}{\mathbf{G}}
\newcommand{\mT}{\mathbf{T}}
\newcommand{\mI}{\mathbf{I}}
\newcommand{\mB}{\mathbf{B}}
\newcommand{\mK}{\mathbf{K}}
\newcommand{\mP}{\mathbf{P}}
\newcommand{\mF}{\mathbf{F}}
\newcommand{\mE}{\mathbf{E}}
\newcommand{\mA}{\mathbf{A}}
\newcommand{\mH}{\mathbf{H}}
% Sets and groups
\newcommand{\RR}{\mathbb{R}}
\newcommand{\SO}{\mathrm{SO}}
% Units
\newcommand{\degps}{\,\text{deg/s}}
\newcommand{\Nm}{\ensuremath{\,\text{N}\cdot\text{m}}}
\newcommand{\BW}{\,\text{BW}}

% ========== Header/Footer ==========
\pagestyle{fancy}
\fancyhf{}
\fancyhead[L]{\leftmark}
\fancyhead[R]{\thepage}
\renewcommand{\headrulewidth}{0.4pt}

% ========== Title formatting ==========
\titleformat{\chapter}[display]
  {\normalfont\huge\bfseries\color{defblue}}
  {\chaptertitlename\ \thechapter}{20pt}{\Huge}

\titleformat{\section}
  {\normalfont\Large\bfseries\color{defblue!80!black}}
  {\thesection}{1em}{}

\titleformat{\subsection}
  {\normalfont\large\bfseries\color{defblue!60!black}}
  {\thesubsection}{1em}{}


\begin{document}

% ==================== TITLE PAGE ====================
\begin{titlepage}
\centering
\vspace*{1.5cm}
{\Large\color{gray} 수학의 정석 시리즈}
\vspace{0.5cm}

{\Huge\bfseries\color{defblue} 다시점 카메라 보정에서\\[0.3cm]
SMPL 인체 피팅까지\\[0.3cm]
수학적 완전 가이드}

\vspace{1cm}

{\LARGE\color{defblue!70!black} 투영 기하학, 좌표 변환, 삼각측량,\\
그리고 2D 재투영 기반 인체 모델 피팅의\\
모든 수식을 빈틈없이 유도합니다}

\vspace{3cm}
{\large 2026}

\vspace{\fill}
{\small\color{gray} 이 문서는 ``왜 이렇게 해야 하는가''를 수식으로 증명합니다.\\
직관과 수식이 함께 흐르는, 수학의 정석 스타일의 완전 해설서입니다.}
\end{titlepage}

\tableofcontents
\newpage


% ====================================================================
\chapter{핀홀 카메라 모델}
\label{ch:pinhole}
% ====================================================================

\section{3D 점에서 2D 픽셀로: 투영의 수학}

\begin{definition}[핀홀 카메라 투영]
3D 월드 좌표 $\vX = (X, Y, Z)^\top \in \RR^3$가 2D 이미지 좌표 $\vu = (u, v)^\top \in \RR^2$로 투영되는 과정은 다음과 같다:

\begin{equation}
\lambda \begin{pmatrix} u \\ v \\ 1 \end{pmatrix}
= \underbrace{\begin{pmatrix} f_x & 0 & c_x \\ 0 & f_y & c_y \\ 0 & 0 & 1 \end{pmatrix}}_{\mK \text{ (내부 행렬)}}
\underbrace{\begin{pmatrix} \mR & \vt \end{pmatrix}}_{\text{외부 행렬 } 3 \times 4}
\begin{pmatrix} X \\ Y \\ Z \\ 1 \end{pmatrix}
\label{eq:projection}
\end{equation}

여기서 $\lambda > 0$는 깊이(depth) 스케일 팩터이다.
\end{definition}

\subsection{내부 파라미터 $\mK$의 의미}

\begin{itemize}[leftmargin=*]
  \item $f_x, f_y$: \textbf{초점 거리} (focal length, 픽셀 단위). 렌즈에서 센서까지의 거리를 센서 픽셀 크기로 나눈 것.
  \begin{equation}
    f_x = \frac{f}{s_x}, \quad f_y = \frac{f}{s_y}
  \end{equation}
  여기서 $f$는 물리적 초점 거리(mm), $s_x, s_y$는 센서 픽셀의 물리적 크기(mm/px).
  정사각 픽셀이면 $f_x = f_y$.

  \item $c_x, c_y$: \textbf{주점} (principal point). 광축이 이미지 평면과 만나는 점.
  이상적으로 이미지 중심 $(W/2, H/2)$이지만, 렌즈 조립 오차로 약간 다를 수 있음.
\end{itemize}

\subsection{외부 파라미터 $[\mR | \vt]$의 의미}

\begin{itemize}[leftmargin=*]
  \item $\mR \in SO(3)$: 월드 좌표계에서 카메라 좌표계로의 \textbf{회전 행렬} ($3 \times 3$, 정규직교, $\det(\mR) = 1$)
  \item $\vt \in \RR^3$: \textbf{이동 벡터}. 카메라 좌표계에서 본 월드 원점의 위치.
\end{itemize}

\begin{importantnote}[title={카메라 위치와 $\vt$의 관계}]
$\vt$는 카메라의 위치가 \textbf{아니다}.
카메라의 월드 좌표계 위치 $\mathbf{C}$는:
\begin{equation}
\mathbf{C} = -\mR^\top \vt
\label{eq:camera_center}
\end{equation}

$[\mR | \vt]$는 ``월드 $\rightarrow$ 카메라'' 변환이므로, $\vt$는 월드 원점을 카메라 좌표로 표현한 것이다.
\end{importantnote}


\section{투영 행렬 $\mP$}

\begin{equation}
\mP = \mK [\mR | \vt] \in \RR^{3 \times 4}
\label{eq:proj_matrix}
\end{equation}

이 행렬 하나로 3D $\rightarrow$ 2D 변환이 완결된다:
\begin{equation}
\lambda \tilde{\vu} = \mP \tilde{\vX}
\end{equation}
여기서 $\tilde{\vu} = (u, v, 1)^\top$, $\tilde{\vX} = (X, Y, Z, 1)^\top$는 동차 좌표(homogeneous coordinates).

\begin{proofbox}[title={투영의 전개}]
$\mP$의 행을 $\mathbf{p}_1^\top, \mathbf{p}_2^\top, \mathbf{p}_3^\top$으로 놓으면:
\begin{align}
u &= \frac{\mathbf{p}_1^\top \tilde{\vX}}{\mathbf{p}_3^\top \tilde{\vX}}, \qquad
v = \frac{\mathbf{p}_2^\top \tilde{\vX}}{\mathbf{p}_3^\top \tilde{\vX}}
\label{eq:proj_explicit}
\end{align}
이것이 비선형 투영 함수 $\pi: \RR^3 \rightarrow \RR^2$를 정의한다.
\end{proofbox}


% ====================================================================
\chapter{이미지 변환과 카메라 파라미터의 변화}
\label{ch:transforms}
% ====================================================================

\begin{keyidea}
이미지를 리사이즈하거나 크롭하면 $\mK$가 변한다.
$[\mR | \vt]$는 이미지 변환에 \textbf{영향을 받지 않는다} --- 이것은 3D 공간의 속성이다.
\end{keyidea}

\section{리사이즈 (Resize)}

\begin{theorem}[리사이즈와 $\mK$의 관계]
\label{thm:resize}
$W \times H$ 이미지를 $W' \times H'$로 리사이즈하면, 새로운 내부 행렬은:
\begin{equation}
\mK' = \begin{pmatrix} s_x & 0 & 0 \\ 0 & s_y & 0 \\ 0 & 0 & 1 \end{pmatrix} \mK
= \begin{pmatrix} s_x f_x & 0 & s_x c_x \\ 0 & s_y f_y & s_y c_y \\ 0 & 0 & 1 \end{pmatrix}
\label{eq:resize_K}
\end{equation}
여기서 $s_x = W'/W$, $s_y = H'/H$.
\end{theorem}

\begin{proofbox}
리사이즈는 2D 좌표에 스케일을 적용하는 것이다:
$u' = s_x u$, $v' = s_y v$.

투영 $\lambda (u, v, 1)^\top = \mK[\mR|\vt]\tilde{\vX}$에서
$\lambda (u', v', 1)^\top = \lambda (s_x u, s_y v, 1)^\top$를 만족하려면,
$\mK$에 좌측에서 대각 스케일 행렬을 곱해야 한다. $\square$
\end{proofbox}

\begin{practicebox}[title={실전: DA3의 리사이즈}]
DA3의 \texttt{upper\_bound\_resize}는 longest side를 \texttt{process\_res}로 맞춘다.

landscape $1920 \times 1080$, \texttt{process\_res}$= 504$:
\begin{align}
s &= \frac{504}{1920} = 0.2625 \\
W' &= \text{round}(1920 \times 0.2625) = 504 \\
H' &= \text{round}(1080 \times 0.2625) = 284 \xrightarrow{\text{div by 14}} 280
\end{align}

이때 $\mK$의 변환:
\begin{equation}
f_x' = f_x \cdot \frac{504}{1920}, \quad
f_y' = f_y \cdot \frac{280}{1080}, \quad
c_x' = c_x \cdot \frac{504}{1920}, \quad
c_y' = c_y \cdot \frac{280}{1080}
\end{equation}

\textbf{주의}: $s_x \neq s_y$ (504/1920 $\neq$ 280/1080)이므로 $f_x' / f_y' \neq f_x / f_y$.
비율 보존 리사이즈에서도 \texttt{make\_divisible\_by\_14}가 높이를 284$\rightarrow$280으로 바꾸면서
$s_y$가 미세하게 달라진다.
\end{practicebox}


\section{크롭 (Crop)}

\begin{theorem}[크롭과 $\mK$의 관계]
\label{thm:crop}
이미지에서 $(x_0, y_0)$을 좌상단으로 하는 $W' \times H'$ 영역을 크롭하면:
\begin{equation}
\mK' = \begin{pmatrix} 1 & 0 & -x_0 \\ 0 & 1 & -y_0 \\ 0 & 0 & 1 \end{pmatrix} \mK
= \begin{pmatrix} f_x & 0 & c_x - x_0 \\ 0 & f_y & c_y - y_0 \\ 0 & 0 & 1 \end{pmatrix}
\label{eq:crop_K}
\end{equation}

\textbf{$f_x, f_y$는 변하지 않는다.} $c_x, c_y$만 이동한다.
\end{theorem}

\begin{proofbox}
크롭은 2D 좌표를 평행 이동하는 것이다: $u' = u - x_0$, $v' = v - y_0$.
$\mK$에 좌측에서 이동 행렬을 곱하면 주점만 변한다. $\square$
\end{proofbox}


\section{Center Crop}

\begin{corollary}[Center crop]
$W \times H$ 이미지를 $W' \times H'$로 center crop하면:
\begin{equation}
x_0 = \frac{W - W'}{2}, \quad y_0 = \frac{H - H'}{2}
\end{equation}
\begin{equation}
c_x' = c_x - \frac{W - W'}{2}, \quad c_y' = c_y - \frac{H - H'}{2}
\label{eq:center_crop}
\end{equation}
\end{corollary}


\section{Letterbox 패딩 (Padding)}

\begin{theorem}[패딩과 $\mK$의 관계]
\label{thm:padding}
이미지 주변에 패딩을 추가하는 것은 크롭의 \textbf{역연산}이다.
$W \times H$ 이미지를 $W' \times H'$로 패딩(중앙 정렬)하면:
\begin{equation}
\text{pad}_x = \frac{W' - W}{2}, \quad \text{pad}_y = \frac{H' - H}{2}
\end{equation}
\begin{equation}
c_x' = c_x + \text{pad}_x, \quad c_y' = c_y + \text{pad}_y
\label{eq:padding_K}
\end{equation}

$f_x, f_y$는 변하지 않는다 (패딩은 리사이즈가 아니므로).
\end{theorem}


\section{복합 변환}

\begin{theorem}[리사이즈 $\rightarrow$ 크롭 복합 변환]
\label{thm:composite}
리사이즈 $(s_x, s_y)$ 후 center crop $(x_0, y_0)$을 적용하면:
\begin{equation}
\mK_{\text{final}} = \mT_{\text{crop}} \cdot \mS_{\text{resize}} \cdot \mK_{\text{orig}}
\end{equation}
\begin{equation}
\boxed{
\begin{aligned}
f_x' &= s_x \cdot f_x \\
f_y' &= s_y \cdot f_y \\
c_x' &= s_x \cdot c_x - x_0 \\
c_y' &= s_y \cdot c_y - y_0
\end{aligned}
}
\label{eq:composite}
\end{equation}
\end{theorem}

\begin{practicebox}[title={실전: DA3 280$\times$280의 역변환}]
DA3가 landscape 1920$\times$1080을 280$\times$280으로 변환한 과정:
\begin{enumerate}[nosep]
  \item 리사이즈: $(s_x, s_y) = (504/1920, 280/1080)$
  \item Center crop: $(x_0, y_0) = ((504-280)/2, 0) = (112, 0)$
\end{enumerate}

역변환 (280$\times$280 $\rightarrow$ 1920$\times$1080):
\begin{enumerate}[nosep]
  \item Uncrop: $c_x \mathrel{+}= 112$, $c_y$ 변경 없음
  \item Unresize: $f_x \mathrel{/}= (504/1920)$, $c_x \mathrel{/}= (504/1920)$, $f_y \mathrel{/}= (280/1080)$, $c_y \mathrel{/}= (280/1080)$
\end{enumerate}

\textbf{이 역변환이 정확해야 EasyMocap의 재투영 피팅이 작동한다.}
\end{practicebox}

\begin{warningbox}
\textbf{DA3의 배치 통일(center crop to min size)} 문제:

DA3가 배치 처리 시 서로 다른 크기의 이미지를 가장 작은 크기로 center crop한다.
이때 landscape(504$\times$280)와 portrait(280$\times$504)가 섞이면
280$\times$280으로 crop되어 \textbf{landscape에서 좌우 224px이 잘린다}.

잘린 영역의 2D 키포인트는 280$\times$280 좌표계에 존재하지 않으므로,
이 카메라 파라미터로는 원본 해상도 키포인트와 호환되지 않는다.

\textbf{해결}: Letterbox 패딩으로 모든 이미지를 같은 정사각형으로 만들어 DA3에 입력.
\end{warningbox}


% ====================================================================
\chapter{Letterbox 패딩과 DA3 보정}
\label{ch:letterbox}
% ====================================================================

\section{Letterbox 전략}

모든 이미지를 $S \times S$ 정사각형으로 만든다 ($S = \max(W_{\max}, H_{\max})$).

\begin{definition}[Letterbox 패딩]
$W \times H$ 이미지를 $S \times S$로 패딩:
\begin{equation}
\text{pad}_x = \frac{S - W}{2}, \quad \text{pad}_y = \frac{S - H}{2}
\end{equation}
패딩 영역은 검은색(0,0,0)으로 채운다.
\end{definition}

현재 세팅에서:
\begin{table}[H]
\centering
\renewcommand{\arraystretch}{1.3}
\begin{tabular}{lcccc}
\toprule
\textbf{카메라} & \textbf{원본} & \textbf{패딩 후} & \textbf{pad$_x$} & \textbf{pad$_y$} \\
\midrule
cam1 (portrait) & 1080$\times$1920 & 1920$\times$1920 & 420 & 0 \\
cam2--7 (landscape) & 1920$\times$1080 & 1920$\times$1920 & 0 & 420 \\
\bottomrule
\end{tabular}
\end{table}


\section{DA3 출력에서 원본 좌표계로의 변환}

DA3가 $S \times S$ 이미지에서 추정한 $\mK_{\text{DA3}}$를 원본 좌표계로 변환:

\begin{equation}
\boxed{
\begin{aligned}
f_x^{\text{orig}} &= f_x^{\text{DA3}} \quad (\text{패딩은 스케일 변경 없음}) \\
f_y^{\text{orig}} &= f_y^{\text{DA3}} \\
c_x^{\text{orig}} &= c_x^{\text{DA3}} - \text{pad}_x \\
c_y^{\text{orig}} &= c_y^{\text{DA3}} - \text{pad}_y
\end{aligned}
}
\label{eq:letterbox_to_orig}
\end{equation}

\textbf{$\mR, \vt$는 변하지 않는다} --- 이것들은 3D 공간의 속성이다.

\begin{keyidea}
\textbf{왜 Letterbox가 올바른 해결책인가}:
\begin{enumerate}[nosep]
  \item 리사이즈가 없으므로 $f_x, f_y$가 \textbf{그대로} 원본 해상도에서 유효
  \item 원본 좌표로의 변환이 \textbf{$c_x, c_y$에서 정수 offset을 빼는 것}만으로 완결
  \item 모든 원본 이미지 영역이 DA3 입력에 포함됨 (잘림 없음)
  \item 배치 내 모든 이미지가 같은 크기이므로 DA3의 center crop이 발생하지 않음
\end{enumerate}
\end{keyidea}


\section{DA3의 내부 처리}

\begin{warningbox}
DA3는 \texttt{make\_divisible\_by\_14} 단계에서 $S \rightarrow S'$ 미세 리사이즈를 수행한다.
$S = 1920$이면 $S' = 1918$ (1918/14 = 137.0).

이 리사이즈의 스케일:
\begin{equation}
s = \frac{S'}{S} = \frac{1918}{1920} = 0.99896
\end{equation}

이 미세한 스케일이 $\mK$에 적용된다:
\begin{equation}
f_x^{\text{DA3}} = s \cdot f_x^{\text{padded}}, \quad c_x^{\text{DA3}} = s \cdot c_x^{\text{padded}}
\end{equation}

따라서 정밀한 역변환:
\begin{equation}
f_x^{\text{orig}} = f_x^{\text{DA3}} / s, \quad c_x^{\text{orig}} = c_x^{\text{DA3}} / s - \text{pad}_x
\end{equation}

다만 $s = 0.999$이므로 실용적으로는 무시 가능한 수준이다 (0.1\% 오차).
\end{warningbox}


% ====================================================================
\chapter{삼각측량 (Triangulation)}
\label{ch:triangulation}
% ====================================================================

\section{DLT 삼각측량}

\begin{definition}[DLT 삼각측량 문제]
$C$개 카메라에서 관찰된 2D 좌표 $\{(u_c, v_c)\}_{c=1}^{C}$와
투영 행렬 $\{\mP_c\}_{c=1}^{C}$가 주어질 때,
3D 점 $\vX$를 추정하는 문제.
\end{definition}

각 카메라 $c$의 투영 관계 $\lambda_c \tilde{\vu}_c = \mP_c \tilde{\vX}$에서
$\lambda_c$를 소거하면:

\begin{equation}
u_c (\mathbf{p}_{3,c}^\top \tilde{\vX}) - \mathbf{p}_{1,c}^\top \tilde{\vX} = 0, \qquad
v_c (\mathbf{p}_{3,c}^\top \tilde{\vX}) - \mathbf{p}_{2,c}^\top \tilde{\vX} = 0
\end{equation}

$C$개 카메라로부터 $2C$개 방정식을 쌓으면:

\begin{equation}
\mA \tilde{\vX} = \mathbf{0}, \quad
\mA \in \RR^{2C \times 4}
\label{eq:dlt}
\end{equation}

\begin{equation}
\mA = \begin{pmatrix}
u_1 \mathbf{p}_{3,1}^\top - \mathbf{p}_{1,1}^\top \\
v_1 \mathbf{p}_{3,1}^\top - \mathbf{p}_{2,1}^\top \\
\vdots \\
u_C \mathbf{p}_{3,C}^\top - \mathbf{p}_{1,C}^\top \\
v_C \mathbf{p}_{3,C}^\top - \mathbf{p}_{2,C}^\top
\end{pmatrix}
\end{equation}

\begin{theorem}[DLT 해]
$\mA$의 SVD $\mA = \mathbf{U} \bm{\Sigma} \mathbf{V}^\top$에서,
$\mathbf{V}$의 마지막 열이 $\tilde{\vX}$의 최소 노름 해이다:
\begin{equation}
\tilde{\vX}^* = \mathbf{v}_4, \quad \vX^* = \frac{1}{v_{4,4}} (v_{4,1}, v_{4,2}, v_{4,3})^\top
\end{equation}
\end{theorem}


\section{재투영 오차}

\begin{definition}[재투영 오차]
삼각측량된 3D 점 $\vX^*$를 각 카메라로 재투영한 2D 좌표와
검출된 2D 좌표의 차이:
\begin{equation}
e_c = \|\pi_c(\vX^*) - \vu_c\|_2 \quad (\text{pixels})
\label{eq:reproj_error}
\end{equation}
\end{definition}

\begin{keyidea}
\textbf{재투영 오차는 카메라 파라미터의 정확도를 직접 측정한다.}

$\mK$나 $[\mR|\vt]$가 틀리면:
\begin{itemize}[nosep]
  \item $\mP_c = \mK_c[\mR_c|\vt_c]$가 틀어짐
  \item DLT 삼각측량이 잘못된 $\vX^*$를 추정
  \item 재투영 시 $\pi_c(\vX^*)$와 검출 $\vu_c$가 불일치
  \item 재투영 오차가 큼
\end{itemize}

\textbf{판정 기준}:
$< 5$\,px (우수), $5$--$20$\,px (수용), $> 50$\,px (카메라 파라미터 불일치)
\end{keyidea}


\section{삼각측량의 조건}

\begin{proposition}[최소 카메라 수]
$\mA \in \RR^{2C \times 4}$의 rank가 최소 3이어야 비자명 해가 존재한다.
$C \geq 2$이면 rank $\leq 4$이므로 해가 존재하지만,
$C = 2$일 때 baseline이 작으면 깊이 추정 오차가 크다.
$C \geq 3$이 실용적 최소 조건이다.
\end{proposition}

\begin{proposition}[7대 카메라의 이점]
$C = 7$이면 $\mA \in \RR^{14 \times 4}$로 과결정(overdetermined).
$\binom{7}{2} = 21$개의 에피폴라 제약이 깊이를 강력하게 결정한다.
\end{proposition}


% ====================================================================
\chapter{2D 재투영 기반 SMPL 피팅}
\label{ch:smpl_fitting}
% ====================================================================

\section{왜 2D 재투영인가}

\begin{theorem}[3D 피팅 vs 2D 재투영 피팅의 방정식 수]
\label{thm:equations}
~\\
\textbf{3D 직접 피팅}: OpenPose 25관절 중 14개 매핑 가능 $\rightarrow$ $14 \times 3 = 42$개 방정식.
미지수: $\vtheta$(72) + $\vbeta$(10) + $\mR$(3) + $\vt$(3) = 88.
$\rightarrow$ \textbf{부정계 (42 $<$ 88)}, prior 필수.

\textbf{다시점 2D 재투영 피팅}: $7 \times 25 \times 2 = 350$개 방정식.
미지수: 88 (동일).
$\rightarrow$ \textbf{과결정 (350 $\gg$ 88)}, prior 없이도 해 결정 가능.
\end{theorem}

\section{EasyMocap의 목적 함수}

\begin{equation}
\boxed{
(\vtheta^*, \vbeta^*) = \arg\min_{\vtheta, \vbeta}
\sum_{c=1}^{7} \sum_{k=1}^{K} w_{ck} \left\| \vu_{ck} - \pi_c\!\big(J_k(\vtheta, \vbeta)\big) \right\|^2
+ \lambda_\theta \|\vtheta\|^2 + \lambda_\beta \|\vbeta\|^2
}
\label{eq:easymocap}
\end{equation}

여기서:
\begin{itemize}[nosep]
  \item $\vu_{ck} \in \RR^2$: 카메라 $c$에서 관절 $k$의 2D 검출 좌표 (OpenPose)
  \item $\pi_c(\cdot) = \mK_c [\mR_c | \vt_c] (\cdot)$: 카메라 $c$의 투영 함수
  \item $J_k(\vtheta, \vbeta) \in \RR^3$: SMPL forward kinematics의 $k$번째 관절 3D 위치
  \item $w_{ck}$: 검출 신뢰도 (OpenPose confidence)
\end{itemize}

\begin{keyidea}
이 수식에서 $\mK_c$가 $\vu_{ck}$와 \textbf{같은 좌표계}여야 하는 이유:

$\pi_c(J_k)$는 $\mK_c$를 사용하여 3D$\rightarrow$2D를 계산한다.
$\vu_{ck}$는 OpenPose가 원본 이미지에서 검출한 좌표이다.
$\mK_c$가 다른 해상도/크롭 기준이면,
$\pi_c(J_k)$와 $\vu_{ck}$가 다른 좌표계에 있게 되어
\textbf{어떤 $(\vtheta, \vbeta)$도 오차를 0으로 만들 수 없다}.

이것이 카메라 파라미터 변환이 정확해야 하는 수학적 이유이다.
\end{keyidea}


\section{3D 직접 피팅이 구조적으로 실패하는 이유}

\begin{theorem}[관절 정의 오프셋]
OpenPose의 관절 $j$의 3D 위치 $\vp_j^{\text{OP}}$와
SMPL의 대응 관절 $J_k(\vtheta, \vbeta)$ 사이에는 체계적 오프셋 $\bm{\delta}_k$가 존재한다:
\begin{equation}
\vp_j^{\text{OP}} = J_k(\vtheta, \vbeta) + \bm{\delta}_k + \bm{\epsilon}
\end{equation}
여기서 $\bm{\delta}_k$는 관절 정의 차이 (시각적 랜드마크 vs 골격 회전 중심),
$\bm{\epsilon}$은 삼각측량 노이즈.

3D 피팅의 목적 함수 $\|J_k - \vp_j^{\text{OP}}\|^2$는 $\bm{\delta}_k$ 때문에 0에 도달할 수 없다.
\end{theorem}

\begin{proofbox}[title={2D 재투영에서 오프셋이 희석되는 이유}]
2D 재투영에서는:
\begin{equation}
\pi_c(J_k + \bm{\delta}_k) - \pi_c(J_k) \approx \mJ_\pi \bm{\delta}_k
\end{equation}
여기서 $\mJ_\pi$는 투영의 야코비안이다.
카메라 $c$마다 $\mJ_\pi$가 다르므로, $\bm{\delta}_k$의 2D 투영이 시점마다 다른 방향으로 나타난다.
7대 카메라에서 이 오프셋들이 \textbf{평균적으로 상쇄}되어
$(\vtheta, \vbeta)$ 추정에 미치는 영향이 줄어든다.

반면 3D에서는 $\bm{\delta}_k$가 모든 프레임에서 동일한 방향으로 누적된다. $\square$
\end{proofbox}


% ====================================================================
\chapter{전체 파이프라인 요약}
\label{ch:pipeline}
% ====================================================================

\begin{figure}[H]
\centering
\begin{tikzpicture}[
  block/.style={rectangle, rounded corners=3pt, draw=defblue, fill=lightblue,
    minimum width=3.5cm, minimum height=0.8cm, align=center, font=\small},
  math/.style={rectangle, rounded corners=3pt, draw=cvpurple, fill=cvpurple!8,
    minimum width=3.5cm, minimum height=0.8cm, align=center, font=\small},
  arrow/.style={-{Stealth[length=2.5mm]}, thick, defblue!70}
]
\node[block] (orig) at (0,0) {원본 이미지\\$W \times H$};
\node[block] (pad) at (0,-1.5) {Letterbox 패딩\\$S \times S$};
\node[block] (da3) at (0,-3) {DA3 추론\\$\mK_{\text{DA3}}, [\mR|\vt]$};
\node[math] (unpad) at (0,-4.5) {Unpad: $c_x \mathrel{-}= \text{pad}_x$\\$c_y \mathrel{-}= \text{pad}_y$};
\node[block] (K_orig) at (0,-6) {$\mK_{\text{orig}}$\\원본 해상도 기준};

\node[block] (op2d) at (6,-1.5) {OpenPose 2D\\$\vu_{ck}$ (원본 좌표)};
\node[math] (proj) at (3,-7.5) {$\pi_c(\cdot) = \mK_{\text{orig}}[\mR_c|\vt_c](\cdot)$};
\node[block] (smpl) at (3,-9) {SMPL 피팅\\$\min \sum \|u_{ck} - \pi_c(J_k)\|^2$};

\draw[arrow] (orig) -- (pad);
\draw[arrow] (pad) -- (da3);
\draw[arrow] (da3) -- (unpad);
\draw[arrow] (unpad) -- (K_orig);
\draw[arrow] (K_orig) -- (proj);
\draw[arrow] (op2d) -- (proj);
\draw[arrow] (proj) -- (smpl);

\node[font=\scriptsize\color{darkred}] at (-2.3, -4.5) {핵심 변환};
\node[font=\scriptsize\color{darkred}] at (4.8, -7.5) {좌표계 일치 필수};
\end{tikzpicture}
\caption{Letterbox DA3 보정 $\rightarrow$ 좌표 변환 $\rightarrow$ EasyMocap SMPL 피팅 전체 흐름}
\end{figure}

\begin{keyidea}
\textbf{전체 파이프라인의 수학적 일관성 조건}:

$\mK_{\text{orig}}$의 좌표계 $=$ $\vu_{ck}$의 좌표계 $=$ 원본 이미지의 픽셀 좌표계

이 세 가지가 일치할 때, 그리고 \textbf{오직 그때만},
EasyMocap의 재투영 오차 $\|\vu_{ck} - \pi_c(J_k)\|$가 의미를 가지며,
SMPL 파라미터가 올바르게 수렴한다.
\end{keyidea}


\end{document}
